%==============================================================================
% Sjabloon poster bachproef
%==============================================================================
% Gebaseerd op document class `a0poster' door Gerlinde Kettl en Matthias Weiser
% Aangepast voor gebruik aan HOGENT door Jens Buysse en Bert Van Vreckem

\documentclass[a0,portrait]{hogent-poster}

% Info over de opleiding
\course{Bachelor thesis}
\studyprogramme{Bachelor of applied computer science}
\academicyear{2025-2026}
\institution{Hogeschool Gent, Valentin Vaerwyckweg 1, 9000 Gent}

% Info over de bachelorproef
\title{To what extent does Meshtastic, as a decentralized LoRa-based mesh network, offer a reliable and readily deployable emergency communication solution when primary network infrastructure fails?}
\author{Simon Dierickx}
\email{simon.dierickx@student.hogent.be}
\supervisor{Steven Van Impe}
\cosupervisor{Moïse Aubert (Accenture)}

% Indien ingevuld, wordt deze informatie toegevoegd aan het einde van de
% abstract. Zet in commentaar als je dit niet wilt.
\specialisation{System \& Network Administrator}
\keywords{Meshtastic, LoRa, Mesh network, Emergency, Off-grid}
\projectrepo{https://github.com/SimonDierickx/bp-2526-DierickxSimon}

\begin{document}

\maketitle

\begin{abstract}
Emergency communication relies heavily on existing network infrastructure, the same infrastructure that fails during storms, power outages, large-scale incidents or deliberate shutdowns.
Existing off-grid alternatives (APRS, DMR, satellite messengers) all fail at least one key requirement: affordability, licence-free use, or easy deployment. This thesis evaluates Meshtastic, a decentralized LoRa-based mesh network, through a real-world three-node proof of concept in the Flemish Ardennes and Pays des Collines.
Across three field sessions the range, RSSI, SNR and packet delivery rate were measured for both direct and relayed links.
\end{abstract}

\begin{multicols}{2} % This is how many columns your poster will be broken into, a portrait poster is generally split into 2 columns

\section{Introduction}
 
When exisiting infrastructure fails, power is cut, or networks are (deliberately) shut down, first responders and communities need a fallback that is independent of fixed infrastructure, easy to deploy, and affordable at scale.
 
Existing off-grid options fall short on important requirements: APRS and DMR require an amateur radio licence and specialized hardware, and satellite messengers carry high per-device and subscription costs.
Meshtastic is an open-source, LoRa-based mesh platform that operates in the licence-free 868\,MHz band, needs no SIM card or internet, and self-heals through multi-hop routing.
Its real-world reliability in an emergency context had not been evaluated in a real-world setting, which is the gap this thesis addresses.

\section{Method}
 
A three-node proof of concept was deployed:
 
\begin{itemize}
  \item A fixed \textbf{base node} in Ronse (90\,m) (autonomous power).
  \item A fixed \textbf{relay node} on the Mont-Saint-Aubert near Tournai (140\,m), mounted with permission in a viewing tower (autonomous power).
  \item A car-mounted \textbf{mobile node}.
\end{itemize}
 
All nodes ran the EU\_868 LONG\_FAST preset. Across three field sessions, range, RSSI, SNR and packet delivery rate (PDR) were measured for both a direct base-to-mobile link and a base-to-relay-to-mobile chain.
The fixed nodes were solar and battery powered for autonomous operation.

\section{Results}
 
\begin{itemize}
  \item A \textbf{direct link of 18.5\,km} was achieved, with at 99.7\,\% PDR and a average SNR of $+6$\,dB.
  \item A single relay extended the \textbf{range to 30\,km}, larger than the literature suggested.
  \item \textbf{Line-of-sight dominates over altitude and distance}: the three highest test sites (up to 140\,m) failed completely due to tree cover, while lower sites with a clear line-of-sight kept a working link (see figure).
  \item The complete setup (4 nodes, antennas, cables, etc.) cost \textbf{€180.95}, far below a single subscription-based satellite messenger.
\end{itemize}
 
\begin{center}
  \captionsetup{type=figure}
  \includegraphics[width=1.0\linewidth]{plot-los-altitude.jpg}
  \captionof{figure}{Altitude and line-of-sight vs. connection (Session~2). The highest sites failed, while lower clear-line-of-sight sites kept a working link.}
\end{center}

\section{Conclusion}
 
Within the constraints, Meshtastic is a credible, low-cost, infrastructure-independent option for emergency communication.
It wins on cost, licence-free operation and deployability requirements, but trades away throughput, voice, and guaranteed delivery.
Reliable coverage depends on line-of-sight and sufficient node placement rather than raw distance or altitude. For a low-cost, text-based network that survives infrastructure failure, this trade-off is viable.
 
\section{Future research}
 
A second relay could extend the chain to an estimated 45--50\,km over three hops, though each hop adds latency and packet loss.
Further work includes more hops, testing across different weather conditions, higher-gain or directional antennas, and a deeper analysis of Meshtastic's security features.

\end{multicols}
\end{document}